\documentclass[10pt,a4paper]{amsart}
\usepackage[margin=19mm]{geometry}
\usepackage{amsmath,amssymb,mathtools}
\usepackage[scr=rsfs,cal=euler]{mathalfa}
\usepackage{xfrac}
\usepackage[unicode,hidelinks,bookmarks=false]{hyperref}
\newcommand{\ie}{i.e.\ }
\newcommand{\cf}{cf.\ }
\newcommand{\resp}{respectively\ }
\newcommand{\Subsection}{\subsection}
\providecommand{\nobreakdash}{\nobreak\hbox{-}}
\setlength{\emergencystretch}{2em}
\multlinegap0pt
\usepackage{amssymb}
\usepackage{enumerate}
\usepackage{tikz}
\newtheorem{prop}{Proposition}
\newcommand{\ctr}[1]{{x}_{#1}}
\newcommand{\dfct}{\mathrm{dfct}}
\newcommand{\inv}{\textsc{inv}}
\newcommand{\mfs}[1]{\mathfrak{S}_{#1}}
\newcommand{\nTksp}{\negthickspace\negthickspace}
\newcommand{\ntksp}{\negthickspace}
\newcommand{\ocnet}[1]{\mathcal F_{#1}}%{\mathcal F^{\circ\bullet}_{#1}}
\newcommand{\sn}{\mathfrak{S}_n}
\newcommand{\type}{\mathrm{type}}
\title{On Kazhdan--Lusztig basis elements having no reversal factorization}
\author{Tommy Parisi, Ben Spahiu, Mark Skandera, Jiayuan Wang}
\date{Source: arXiv:2605.21733v1 (CC BY 4.0)}
\begin{document}
\maketitle

\noindent\emph{Source and licence.} On Kazhdan--Lusztig basis elements having no reversal factorization. Version arXiv:2605.21733v1 (CC BY 4.0), distributed by arXiv under the Creative Commons Attribution 4.0 International licence, http://creativecommons.org/licenses/by/4.0/, as recorded in the arXiv licence metadata for this exact version and shown on the abstract page https://arxiv.org/abs/2605.21733v1. Full licence notice: \texttt{licenses/arxiv-2605.21733-CC-BY-4.0.txt}.\par\smallskip

\noindent\textit{Editorial scope.} Counted unit: the article's prop p:qfact (source0.tex lines 2178--2196). The article's complete original proof of it is delivered with it (source0.tex lines 2197--2311, one \texttt{proof} environment of 115 lines). No mathematics is written, repaired or re-derived: the statement and the proof are the article's own lines, in the article's own notation, and no retained line is substituted or re-flowed.  No further source-local context is needed: every cross-reference of every retained line resolves inside the two delivered spans.  Nothing was omitted from any retained span, so the package carries no omission record.  No retained line contains a citation, so the wrapper carries no bibliography and no bibliography entry.  No retained line contains a figure environment, an image-inclusion command or drawing code, so no image is delivered and the package declares no asset dependency.  Editorial formatting only, in addition: an AMS \texttt{amsart} preamble that restates the article's own declarations and macros used by the retained lines, namely the environments \texttt{prop} on separate counters, together with the standard packages those lines need.  The retained environments print in this document as Proposition 1 in order of appearance, whereas the article prints the same items with the numbering of its own full document.  The complete native source of the article as distributed in the arXiv e-print for this exact version is bundled unchanged as \texttt{sources/201080/source0.tex}.
\begin{prop}\label{p:qfact}
    Fix $F \in \ocnet n$, assume that %pair of internal vertices 
    %$\ctr i$, $\ctr j$ 
    %are incident upon 
    %$m_{i, j}$ 
$p > 1$
multiplicity-$1$ edges are incident upon some pair of internal vertices, and construct $F' \in \ocnet n$ by replacing these edges by one edge of multiplicity $p$.
%$m_{i,j}$. 
Then $F'$ graphically represents the element
%    \begin{equation*}
%        \frac 1 {[m_{i, j}]_q!} \sum_{v \in \sn} \sum _{d \geq 0} |\Pi_{v, d}(F)| q^d T_v.
%    \end{equation*}
%    or
    \begin{equation*}
     \frac 1{[p]_q!} 
     %{[m_{i,j}]_q!} 
     \ntksp\sum_{\pi \in \Pi(F)}\ntksp q^{\dfct(\pi)}T_{\type(\pi)}.
     \end{equation*}
\end{prop}

\begin{proof}
    Let $\ctr r$, $\ctr s$ be the pair of internal vertices in $F$, let $e_1, \dots, e_p$ be the edges incident upon both vertices
    %$\ctr r$, $\ctr s$ 
    in $F$, and let $e$ be the multiplicity-$p$ edge replacing these in $F'$.
    %$\ctr i$ and $\ctr j$, and suppose that these are replaced by a single multiplicity-$p$ edge $e$ in creating $F'$. 
    %Observe that 
    Thus for every path family 
    %$\pi \in 
    in $\Pi(F)$,
    $p$ of its component paths contain the edges $e_1,\dotsc, e_p$,
    and for every path family in $\Pi(F')$,
    $p$ of its component paths contain the edge $e$.
    Define the map 
%    \begin{equation*}
%    \begin{aligned}
        $\theta : \Pi(F) \rightarrow \Pi(F')$
        %\\
%        \pi &\mapsto \pi'
%        \end{aligned}
%        \end{equation}
        by replacing edges $e_1,\dotsc, e_p$ in each path family $\pi \in \Pi(F)$ by $p$ copies of the edge $e$.
        %substituting $e$ for these $p$ edges in the $p$ paths.
        
%        and that these preimages partition $\Pi(F)$,
%        \begin{equation*}           
%        \end{equation*}
             It is clear that 
             the map $\theta$ is surjective and that 
             for each path family $\sigma \in \Pi(F')$, 
        %we have
        %\begin{enumerate}
        %that 
        the set $\theta^{-1}(\sigma) \subseteq \Pi(F)$ consists of $p!$ path families $\pi$ satisfying $\type(\pi) =  \type(\sigma)$.
        We claim that the defects of
            path families in each preimage $\theta^{-1}(\sigma)$ satisfy
    \begin{equation}\label{eq:preimage}
        \sum_{\pi \in \theta^{-1}(\sigma)} 
        \nTksp q^{\dfct(\pi)} = [p]_q! q^{\dfct(\sigma)} .
    \end{equation}
    To see this, fix $\pi \in \theta^{-1}(\sigma)$ and observe that for all internal vertices $\ctr t \neq \ctr s$ of $F$ and $F'$, the triple $(\pi_i, \pi_j,  t)$ is defective if and only if $(\sigma_i, \sigma_j, t)$ is defective:
    the agreement of $\pi$ and $\sigma$ outside of the subgraph induced by $\ctr r$, $\ctr s$ guarantees that
    the preorder $\prec_t$ defined in terms of any internal vertex $\ctr t$ satisfies
    $\pi_i \prec_t \pi_j$ if and only if $\sigma_i \prec_t \sigma_j$. It follows that the difference between $\dfct(\pi)$ and
    %is equal to 
    $\dfct(\sigma)$ is just the number of defects of $\pi$ at $\ctr s$.
    Letting $\pi_{u_1}, \dotsc, \pi_{u_p}$ be the $p$ paths passing through $\ctr r$, $\ctr s$, labeled by
%    But this is equal to the number of inversions in the word $u_1 \cdots u_p$ defined by containing $x$ and $y$, with
    %in terms of the preorder $\prec_y$ 
    %defined in terms of the internal vertex $y$ 
    \begin{equation*}
        \pi_{u_1} \prec_s \cdots \prec_s \pi_{u_p},
    \end{equation*}
%Thus 
we have that the number of defects of $\pi$ at $\ctr s$ 
equals the number of inversions in the word $u_1 \cdots u_p$.
Thus we have
\begin{equation*}
\sum_{\pi \in \theta^{-1}(\sigma)} \nTksp q^{\dfct(\pi)} = 
\sum_{u \in \mfs p} q^{\inv(u) + \dfct(\sigma)},
\end{equation*}
which implies \eqref{eq:preimage}.
%as desired.

%By definition and by \eqref{eq:preimage}, 
It follows that
$F'$ graphically represents the element
    \begin{equation}\label{eq:qfactsum}
        % \sum_{v \in \sn} \sum_{d \geq 0} |\Pi_{v,d}(F')|q^d T_v.
        \nTksp\sum_{\sigma \in \Pi(F')}\nTksp q^{\dfct(\sigma)}T_{\type(\sigma)}
    %\end{equation}
   % By \eqref{eq:preimage}, this is 
   % \begin{equation*}
        = \ntksp \sum_{\sigma \in \Pi(F')} 
        %\frac 1 {[m_{i,j}]_q!} 
        \frac 1{[p]_q!}
        \ntksp \sum_{\pi \in \theta^{-1}(\sigma)} \nTksp q^{\dfct(\pi)} T_{\type(\sigma)}.
        \end{equation}
        Since  $\type(\pi) = \type(\sigma)$
        for all $\pi \in \theta^{-1}(\sigma)$, and since
       \begin{equation*}
            \Pi(F) = \nTksp \bigcup_{\sigma \in \Pi(F')}\nTksp \theta^{-1}(\sigma),
        \end{equation*}
        we have the desired result.
%        It is clear that the map $\theta$ is surjective and that for each path family $\sigma \in \Pi(F')$, 
        %we have
        %\begin{enumerate}
        %that 
%        the set $\theta^{-1}(\sigma) \subseteq \Pi(F)$ consists of $p!$ path families $\pi$ satisfying $\type(\pi) =  \type(\sigma)$.
%Thus we have
%        \begin{equation*}
%            \Pi(F) = \nTksp \bigcup_{\sigma \in \Pi(F')}\nTksp \theta^{-1}(\sigma).
%            \end{equation*}
%            \begin{equation}
%                
%            \end{equation}
%            \item for all $\pi \in \theta^{-1}(\sigma)$, $\type(\pi) = \type(\sigma)$,
 %           \item $ \subset \Pi(F)$ contains $p!$ path families of type $\type(\sigma)$.
% Choose $v$ such that $\Pi_v(F)$ is nonempty.
% For each path family $\pi$ in $\Pi_{v}(F)$,
% the $m_{i,j}$ $x_i$-to-$x_j$ edges appear
% in $m_{i,j}$ different paths.
%     Define a map 
%     \begin{equation*}
%     \begin{aligned}
%         \Pi_{v}(F) &\rightarrow \Pi_{v}(F')\\
%         \pi &\mapsto \pi'
%         \end{aligned}
%         \end{equation*}
% by replacing the $m_{i,j}$ edges in the $m_{i,j}$ paths with
% %from $x_i$ to $x_j$ in $m_{i,j}$ paths of $\pi$ with 
% the single $x_i$-to-$x_j$ path in $F'$.
% subpaths of paths in $\pi$ 
%     with $m_{i,j}$ copies
%     $\pi = (\pi_1,\dotsc,\pi_n)$ be a path family covering $F$
\end{proof}

\end{document}
